\documentclass{beamer}

\usepackage{amsmath,amssymb,bm}
\usepackage{physics}
\usepackage{braket}
\usepackage{mathtools}
\usepackage{amsthm}

\usetheme{Madrid}

\title{Unitary Operators, the Spectral Theorem, and HHL}
\subtitle{Eigenphases, Finite/Infinite-Dimensional Structure, and Quantum Linear Systems}
\author{}
\date{}

\begin{document}

\frame{\titlepage}

%================================================
\section{Why \(U\ket{\psi}=e^{2\pi i\phi}\ket{\psi}\)?}
%================================================

\begin{frame}{The Precise Statement}
For a unitary operator \(U\), one often writes
\[
U\ket{\psi}=e^{2\pi i\phi}\ket{\psi}.
\]

This is \emph{not} true for an arbitrary state \(\ket{\psi}\).

\medskip

It is true precisely when \(\ket{\psi}\) is an eigenstate of \(U\).
\end{frame}

\begin{frame}{Eigenvalue Equation}
If \(\ket{\psi}\) is an eigenstate of \(U\), then
\[
U\ket{\psi}=\lambda \ket{\psi}.
\]

For a unitary operator, the eigenvalue must satisfy
\[
|\lambda|=1.
\]

Hence
\[
\lambda=e^{i\theta}
\]
for some real \(\theta\).
\end{frame}

\begin{frame}{Proof that \(|\lambda|=1\)}
Because \(U\) is unitary,
\[
U^\dagger U = I.
\]

If
\[
U\ket{\psi}=\lambda\ket{\psi},
\]
then
\[
\|U\ket{\psi}\|^2 = \|\ket{\psi}\|^2.
\]

But also
\[
\|U\ket{\psi}\|^2 = \|\lambda \ket{\psi}\|^2 = |\lambda|^2\|\ket{\psi}\|^2.
\]

Therefore
\[
|\lambda|^2\|\ket{\psi}\|^2=\|\ket{\psi}\|^2
\quad\Rightarrow\quad
|\lambda|=1.
\]
\end{frame}

\begin{frame}{Why Write \(e^{2\pi i\phi}\)?}
Since \(|\lambda|=1\), every eigenvalue lies on the unit circle:
\[
\lambda=e^{i\theta}.
\]

Because the phase is periodic modulo \(2\pi\),
\[
e^{i(\theta+2\pi)}=e^{i\theta},
\]
we may write
\[
\theta=2\pi\phi,
\qquad
\phi\in[0,1).
\]

Thus
\[
\lambda=e^{2\pi i\phi}.
\]
\end{frame}

\begin{frame}{What About a General State?}
Let \(\{\ket{u_j}\}\) be an eigenbasis of \(U\). For a general state,
\[
\ket{\psi}=\sum_j c_j \ket{u_j}.
\]

Then
\[
U\ket{\psi}=\sum_j c_j e^{2\pi i\phi_j}\ket{u_j}.
\]

So one cannot factor out a single phase unless all occupied eigenspaces carry the same eigenvalue.
\end{frame}

%================================================
\section{Spectral Theorem: Finite Dimension}
%================================================

\begin{frame}{Spectral Theorem in Finite Dimensions}
If \(U\) is unitary on a finite-dimensional Hilbert space, then \(U\) is unitarily diagonalizable.

\medskip

Equivalently, there exists an orthonormal basis \(\{\ket{u_j}\}\) such that
\[
U\ket{u_j}=e^{i\theta_j}\ket{u_j}.
\]

Hence
\[
U=\sum_j e^{i\theta_j}\ket{u_j}\bra{u_j}
=\sum_j e^{2\pi i\phi_j}\ket{u_j}\bra{u_j}.
\]
\end{frame}

\begin{frame}{Why This Holds}
A unitary operator is a normal operator:
\[
U^\dagger U = UU^\dagger.
\]

Every normal operator on a finite-dimensional Hilbert space is unitarily diagonalizable.

Therefore there exists a unitary matrix \(V\) such that
\[
U=VDV^\dagger,
\]
with
\[
D=\mathrm{diag}(e^{i\theta_1},\dots,e^{i\theta_n}).
\]
\end{frame}

\begin{frame}{Projector Form}
The same result may be written as
\[
U=\sum_j \lambda_j P_j,
\]
where
\[
\lambda_j=e^{i\theta_j},
\qquad
P_j=\ket{u_j}\bra{u_j}
\]
in the nondegenerate case.

\medskip

More generally, \(P_j\) projects onto the eigenspace corresponding to \(\lambda_j\).
\end{frame}

\begin{frame}{Finite-Dimensional Proof Sketch}
The finite-dimensional spectral theorem can be organized as follows:
\begin{enumerate}
\item Show that \(U\) is normal.
\item Use Schur decomposition:
\[
U=QTQ^\dagger,
\]
with \(Q\) unitary and \(T\) upper triangular.
\item For a normal operator, the Schur form \(T\) must actually be diagonal.
\item The diagonal entries are the eigenvalues of \(U\), and for a unitary they satisfy \(|\lambda_j|=1\).
\end{enumerate}
\end{frame}

%================================================
\section{Spectral Theorem: Infinite Dimension}
%================================================

\begin{frame}{What Changes in Infinite Dimensions?}
In infinite-dimensional Hilbert spaces, not every unitary operator has a basis of honest eigenvectors.

\medskip

Instead, the spectral theorem is formulated in terms of a projection-valued measure:
\[
U=\int_{\mathbb{T}} z\, dE(z),
\]
where \(\mathbb{T}=\{z\in\mathbb{C}:|z|=1\}\) is the unit circle.
\end{frame}

\begin{frame}{Projection-Valued Measure Form}
The operator-valued measure \(E(\cdot)\) assigns an orthogonal projector to measurable subsets of the unit circle.

Then
\[
U=\int_{\mathbb{T}} z\, dE(z)
\]
plays the role of the discrete sum
\[
U=\sum_j \lambda_j P_j.
\]

Discrete eigenvalues correspond to point masses of the spectral measure, while continuous spectrum is described by the continuous part of \(E\).
\end{frame}

\begin{frame}{Interpretation}
Finite dimension:
\[
U=\sum_j \lambda_j P_j.
\]

Infinite dimension:
\[
U=\int_{\mathbb{T}} z\, dE(z).
\]

So the spectral theorem remains true, but the discrete list of eigenvalues may be replaced by a full spectral distribution on the unit circle.
\end{frame}

\begin{frame}{Example: Translation Operator}
Consider a shift or translation operator on \(L^2\) or \(\ell^2\).

Such operators are unitary, but may fail to possess a normalizable eigenbasis in the usual sense.

\medskip

Their spectral content is still encoded by the projection-valued measure, and Fourier transform often provides the correct diagonal representation.
\end{frame}

\begin{frame}{Why This Matters Physically}
In quantum mechanics:
\begin{itemize}
\item finite-dimensional systems often appear in qubit models,
\item infinite-dimensional systems arise for position, momentum, fields, and wave equations.
\end{itemize}

The same principle survives:
\[
\text{unitary operators are multiplication by phases in the appropriate spectral representation.}
\]
\end{frame}

%================================================
\section{Connection to Time Evolution}
%================================================

\begin{frame}{Unitary Time Evolution}
Time evolution is generated by a Hermitian Hamiltonian \(H\):
\[
U(t)=e^{-iHt}.
\]

If
\[
H\ket{u_j}=E_j\ket{u_j},
\]
then
\[
U(t)\ket{u_j}=e^{-iE_j t}\ket{u_j}.
\]
\end{frame}

\begin{frame}{Spectral Representation of \(U(t)\)}
If
\[
H=\sum_j E_j \ket{u_j}\bra{u_j},
\]
then
\[
U(t)=e^{-iHt}
=\sum_j e^{-iE_j t}\ket{u_j}\bra{u_j}.
\]

In infinite dimensions this becomes
\[
U(t)=\int e^{-iEt}\, dE_H(E),
\]
where \(E_H\) is the spectral measure of \(H\).
\end{frame}

\begin{frame}{Physical Meaning of the Phase}
The quantity
\[
e^{-iE_j t}
\]
is the dynamical phase accumulated by the eigenstate \(\ket{u_j}\).

Thus the eigenphases of the unitary encode the energy spectrum of the generator.
\end{frame}

%================================================
\section{Quantum Phase Estimation}
%================================================

\begin{frame}{Quantum Phase Estimation (QPE)}
QPE begins with the relation
\[
U\ket{u_j}=e^{2\pi i\phi_j}\ket{u_j}.
\]

Its goal is to estimate \(\phi_j\).
\end{frame}

\begin{frame}{Why QPE Works}
The ingredients are:
\begin{itemize}
\item an eigenstate \(\ket{u_j}\),
\item controlled powers \(U^{2^k}\),
\item an inverse quantum Fourier transform.
\end{itemize}

Because
\[
U^{2^k}\ket{u_j}=e^{2\pi i\, 2^k\phi_j}\ket{u_j},
\]
the phase is imprinted on the control register in binary form.
\end{frame}

\begin{frame}{If the Input is Not an Eigenstate}
If
\[
\ket{\psi}=\sum_j c_j \ket{u_j},
\]
then QPE produces
\[
\sum_j c_j \ket{\widetilde{\phi_j}}\ket{u_j},
\]
so the measurement outcome samples one of the eigenphases with probability \(|c_j|^2\).
\end{frame}

%================================================
\section{HHL Algorithm}
%================================================

\begin{frame}{The Linear-System Problem}
We want to solve
\[
A\ket{x}=\ket{b},
\]
where \(A\) is Hermitian and invertible.

Write
\[
A=\sum_j \lambda_j \ket{u_j}\bra{u_j},
\qquad
\ket{b}=\sum_j \beta_j \ket{u_j}.
\]

Then formally
\[
\ket{x}\propto A^{-1}\ket{b}
=\sum_j \frac{\beta_j}{\lambda_j}\ket{u_j}.
\]
\end{frame}

\begin{frame}{The Core Unitary in HHL}
HHL uses the unitary
\[
U=e^{iAt}.
\]

Because \(A\) is Hermitian,
\[
e^{iAt}\ket{u_j}=e^{i\lambda_j t}\ket{u_j}.
\]

Thus the eigenphases are
\[
\phi_j=\frac{\lambda_j t}{2\pi}.
\]
\end{frame}

\begin{frame}{Step 1: State Preparation}
Start with
\[
\ket{b}=\sum_j \beta_j \ket{u_j}.
\]

Attach a clock register and ancilla:
\[
\ket{b}\ket{0}\ket{0}.
\]
\end{frame}

\begin{frame}{Step 2: Phase Estimation}
Apply quantum phase estimation using \(U=e^{iAt}\):
\[
\sum_j \beta_j \ket{u_j}\ket{0}\ket{0}
\;\longrightarrow\;
\sum_j \beta_j \ket{u_j}\ket{\widetilde{\lambda_j}}\ket{0}.
\]

The eigenvalue information is now stored in the second register.
\end{frame}

\begin{frame}{Step 3: Controlled Rotation}
Apply a controlled ancilla rotation so that
\[
\ket{\widetilde{\lambda_j}}\ket{0}
\longrightarrow
\ket{\widetilde{\lambda_j}}
\left(
\sqrt{1-\frac{C^2}{\lambda_j^2}}\ket{0}
+
\frac{C}{\lambda_j}\ket{1}
\right),
\]
where \(C\) is chosen so that \(|C/\lambda_j|\le 1\).
\end{frame}

\begin{frame}{Step 4: Uncomputation}
Undo the phase estimation register:
\[
\sum_j \beta_j \ket{u_j}\ket{\widetilde{\lambda_j}}
\left(
\sqrt{1-\frac{C^2}{\lambda_j^2}}\ket{0}
+
\frac{C}{\lambda_j}\ket{1}
\right)
\]
becomes
\[
\sum_j \beta_j \ket{u_j}\ket{0}
\left(
\sqrt{1-\frac{C^2}{\lambda_j^2}}\ket{0}
+
\frac{C}{\lambda_j}\ket{1}
\right).
\]
\end{frame}

\begin{frame}{Step 5: Postselection}
Condition on measuring the ancilla in the state \(\ket{1}\). Then the system register collapses to
\[
\sum_j \frac{\beta_j}{\lambda_j}\ket{u_j}
\propto A^{-1}\ket{b}.
\]

This is the HHL solution state.
\end{frame}

\begin{frame}{Why the Spectral Theorem is Essential}
Everything in HHL depends on the spectral decomposition
\[
A=\sum_j \lambda_j \ket{u_j}\bra{u_j}.
\]

Without it:
\begin{itemize}
\item QPE would not isolate eigenvalues,
\item controlled rotations would not encode \(1/\lambda_j\),
\item the inverse operator \(A^{-1}\) would not emerge spectrally.
\end{itemize}
\end{frame}

\begin{frame}{Interpretation as a Spectral Filter}
HHL implements the map
\[
\lambda_j \longmapsto \frac{1}{\lambda_j}.
\]

So it is best understood as a spectral filtering algorithm:
\[
A^{-1}=\sum_j \frac{1}{\lambda_j}\ket{u_j}\bra{u_j}.
\]
\end{frame}

\begin{frame}{Connection to Green's Functions}
The inverse operator is analogous to a resolvent or Green's operator:
\[
(\omega I-H)^{-1}.
\]

Thus HHL is closely related to:
\begin{itemize}
\item linear response,
\item propagators,
\item Green's-function methods,
\item discretized differential equations.
\end{itemize}
\end{frame}

\begin{frame}{HHL and Differential Equations}
If a differential equation is discretized into
\[
A\bm{u}=\bm{f},
\]
then HHL prepares a state proportional to the solution vector:
\[
\ket{u}\propto A^{-1}\ket{f}.
\]

The algorithm therefore computes a quantum representation of the discretized solution.
\end{frame}

%================================================
\section{Summary}
%================================================

\begin{frame}{Main Points}
\begin{itemize}
\item \(U\ket{\psi}=e^{2\pi i\phi}\ket{\psi}\) only when \(\ket{\psi}\) is an eigenstate.
\item Eigenvalues of a unitary lie on the unit circle.
\item In finite dimensions:
\[
U=\sum_j e^{2\pi i\phi_j}\ket{u_j}\bra{u_j}.
\]
\item In infinite dimensions:
\[
U=\int_{\mathbb{T}} z\, dE(z).
\]
\item These facts underlie time evolution, QPE, and the HHL algorithm.
\end{itemize}
\end{frame}

\begin{frame}{Conceptual Synthesis}
\[
\boxed{
\text{Spectral theorem} \;\Longrightarrow\;
\text{eigenphases of unitaries} \;\Longrightarrow\;
\text{QPE} \;\Longrightarrow\; \text{HHL}.
}
\]

This is the mathematical chain connecting unitary operators to quantum linear-system solvers.
\end{frame}

\end{document}
